\begin{center}
    Disclaimer
\end{center}

\noindent
I do not guarantee this summary to contain every part of the lecture, as it is only supposed to summarize the things I personally found the most important when going through the slides.

\begin{center}
    This document was compiled on \today.
\end{center}

\newpage

\epigraph{Don't dance faster than the music!}{}
\section*{Introduction}
The Embedded Real-Time Systems lecture of Kiel University by Reinhard von Hanxleden is based on \emph{Introduction to Embedded Systems---A Cyber-Physical Systems Approach} by Edward Lee and Sanjit Seshia\footnote{\url{https://ptolemy.berkeley.edu/books/leeseshia/releases/LeeSeshia_DigitalV2_3.pdf}}.
Embedded systems are computers whose main purpose is the interaction with a physical process.
This interaction is commonly implemented using sensors, actuators, and physical dynamics.
A real-time system is broadly defined by the notion, that not only the correctness of the output of a machine matters, but also the timing in which they arrive.

In this book and lecture it is postulated, that embedded systems should be created using an iterative process, which is called \emph{Modeling, Design, Analysis}.
\begin{description}
    \item[Modeling] The process of gaining a deeper understanding of a system through environmental imitation.
        Models specify \emph{what} a system does.
    \item[Design] The structured creation of artifacts.
        Designs specify \emph{how} a system does what it does.
    \item[Analysis] The process of gaining a deeper understanding of a system through dissection.
        Analysis specifies \emph{why} a system does what it does, or fails to do so.
\end{description}
This approach has the important benefit, that we can implement mechanisms to compile embedded systems software based on a mathematical model of the system and environment.
The compiler can perform a model analysis during translation.
